\documentclass[10pt,letterpaper]{article}
\usepackage[T1]{fontenc}
\usepackage[utf8]{inputenc}
\usepackage[default]{sourcesanspro}
\usepackage{geometry,xcolor,enumitem,titlesec,fancyhdr}
\usepackage[hidelinks]{hyperref}
\geometry{left=0.68in,right=0.68in,top=0.57in,bottom=0.58in}
\definecolor{ink}{HTML}{172A32}
\definecolor{accent}{HTML}{244E5B}
\definecolor{muted}{HTML}{52616A}
\definecolor{rulecolor}{HTML}{D6DFE2}
\hypersetup{pdftitle={Haoyang Jiang - Curriculum Vitae},pdfauthor={Haoyang Jiang},pdfsubject={Scientific machine learning, dynamical systems and inverse problems}}
\setlength{\parindent}{0pt}
\setlength{\parskip}{0pt}
\linespread{1.045}
\pagestyle{fancy}\fancyhf{}\renewcommand{\headrulewidth}{0pt}
\fancyfoot[L]{\small\color{muted}Haoyang Jiang \enspace / \enspace Curriculum Vitae}
\fancyfoot[R]{\small\color{muted}\thepage\,/\,2}
\setlength{\footskip}{22pt}
\titleformat{\section}{\large\bfseries\color{accent}}{}{0pt}{}[\vspace{-3pt}\color{rulecolor}\titlerule]
\titlespacing*{\section}{0pt}{10pt}{6pt}
\setlist[itemize]{leftmargin=12pt,label=\textbullet,itemsep=2pt,topsep=3pt,parsep=0pt,partopsep=0pt}
\newcommand{\role}[2]{\textbf{#1}\hfill{\small\color{muted}#2}\par}
\newcommand{\project}[2]{\vspace{5pt}\textbf{#1}\hfill{\small #2}\par}
\newcommand{\plink}[2]{\href{#1}{\color{accent}#2}}
\newcommand{\me}{\textbf{H. Jiang}}
\begin{document}\color{ink}\raggedright
\hyphenpenalty=10000\exhyphenpenalty=10000
{\fontsize{27}{29}\selectfont\bfseries Haoyang Jiang}\par
\vspace{4pt}
{\large\color{accent}Scientific Machine Learning \enspace | \enspace Dynamical Systems \enspace | \enspace Inverse Problems}\par
\vspace{7pt}
William \& Mary, Williamsburg, Virginia \enspace\textbar\enspace \href{mailto:hjiang16@wm.edu}{hjiang16@wm.edu}\par
\vspace{2pt}
\plink{https://haoyangjiang-wm.github.io/}{Website}\enspace / \enspace
\plink{https://github.com/HaoyangJiang-WM}{GitHub}\enspace / \enspace
\plink{https://scholar.google.com/citations?user=q6923RoAAAAJ}{Google Scholar}\enspace / \enspace
\plink{https://www.linkedin.com/in/haoyang-jiang-7573872a4/}{LinkedIn}\enspace / \enspace
\plink{https://haoyangjiang-wm.github.io/neural-equation-demos/}{Interactive research demos}

\section*{Research Profile}
Ph.D. researcher developing machine-learning methods for physical dynamical systems. Research spans neural operators, physics-guided graph learning, generative inverse modeling, and data assimilation, with applications in hydrology, geophysics, and space systems.

\section*{Education}
\role{William \& Mary}{Expected 2027}
Ph.D. in Data Science \hfill Williamsburg, Virginia\par
\vspace{5pt}
\role{University of Science and Technology of China}{Sep. 2020 -- Jun. 2023}
M.S. in Astrometry and Celestial Mechanics \hfill Hefei, China\par
{\small\color{muted}Research: data assimilation, atmospheric-density uncertainty propagation, and numerical dynamics.}\par
\vspace{5pt}
\role{China University of Mining and Technology (Beijing)}{Sep. 2016 -- Jun. 2020}
B.S. in Geophysics; Minor in Computer Science \hfill Beijing, China\par
{\small\color{muted}Research background: seismic forward modeling, inversion, and computational geophysics.}

\section*{Research Experience}
\project{Physics-guided forecasting on open physical networks}{\plink{https://github.com/HaoyangJiang-WM/OpenBoundary}{Code}}
\begin{itemize}
\item Develop graph-based models for river and coastal systems, combining learned boundary forcing with physics-guided transport and long-horizon forecasting.
\item Study the Danube river network and Chesapeake Bay mesh; the ICDM study reports a 38.3\% average RMSE reduction relative to its spatiotemporal GNN backbone across reported horizons.
\end{itemize}
\project{Neural operators and boundary-value problems}{\plink{https://github.com/HaoyangJiang-WM/CV/blob/main/RESEARCH_PROFILE.md}{Research}}
\begin{itemize}
\item Develop Fredholm-integral-equation neural operators for boundary-value problems and learned ghost-node representations for missing boundary context in physical networks.
\item Connect integral formulations, graph topology, and numerical propagation to improve learning when the modeled domain is incomplete.
\end{itemize}
\project{Generative models and scientific inverse problems}{\plink{https://github.com/HaoyangJiang-WM/FWI}{FWI}\enspace / \enspace\plink{https://github.com/HaoyangJiang-WM/2otcfm}{Flow matching}}
\begin{itemize}
\item Combine pretrained generative priors with inference-time optimization for full-waveform inversion, without paired retraining for each new observation.
\item Investigate second-order optimal-transport flow matching and variational formulations linking position and velocity in generative dynamics.
\end{itemize}
\project{Neural differential, delay, and integral equations}{\plink{https://haoyangjiang-wm.github.io/neural-equation-demos/}{Live demos}}
\begin{itemize}
\item Implement end-to-end Neural ODE, learned-delay DDE, and neural integral models for continuous-time prediction and boundary-driven sensor responses.
\item Evaluate long-horizon and frequency-shift behavior with discrete and implicit adjoints, differentiable history interpolation, gradient checks, and numerical-refinement tests.
\end{itemize}
\project{Satellite dynamics and uncertainty propagation}{}
\begin{itemize}
\item Studied satellite trajectory modeling and orbital-error propagation under atmospheric-density uncertainty using numerical and machine-learning methods.
\end{itemize}

\newpage
\section*{Publications}
\begin{enumerate}[leftmargin=16pt,label={\arabic*.},itemsep=8pt,topsep=1pt,parsep=0pt]
\item \me, Z. Wang, S. Gao, Y. J. Zhang, X. Zhu, and Y. He.
\plink{https://github.com/HaoyangJiang-WM/OpenBoundary/blob/main/ICDM_RefinedGNN.pdf}{\textit{Physics-Refined Spatiotemporal Forecasting on Open-Boundary Hydrologic Graphs.}}
\textbf{IEEE ICDM}, 2026, accepted.\\
{\small\color{accent}Selected among ICDM's best-ranked papers; invited to submit an extended version to KAIS.}
\item \me, B. Qu, X. Zhu, J. Tan, and Y. He.
\plink{https://openreview.net/forum?id=31gTIfhoH0}{\textit{Boundary-Consistent Graph Neural Networks for Topological Flux Prediction.}}
\textbf{Transactions on Machine Learning Research}, 2026.
\item \me\ and Y. Z. Qu.
\textit{Fredholm Integral Equations Neural Operator (FIE-NO) for Boundary Value Problems.}
\textbf{Machine Learning: Engineering}, 2026, accepted.
\item \me, J. D. Wang, Y. He, et al.
\plink{https://proceedings.mlr.press/v267/jiang25i.html}{\textit{Topology-aware Neural Flux Prediction Guided by Physics.}}
\textbf{ICML}, 2025.
\item \me, M. J. Zhang, et al.
\textit{Orbital error propagation considering atmospheric density uncertainty.}
\textbf{Advances in Space Research}, 71(6), 2566--2574, 2023.
\end{enumerate}

\section*{Industry Experience}
\role{Origin Space \enspace | \enspace Assistant Engineer Intern}{Feb. 2022 -- Jun. 2022}
Nanjing, China
\begin{itemize}
\item Applied machine-learning and numerical methods to space-technology research, including reproducing and benchmarking algorithms on engineering data.
\item Improved research prototypes through data preprocessing, parameter tuning, and runtime optimization.
\end{itemize}

\section*{Technical Skills}
\textbf{Machine learning:} Physics-guided learning, neural operators, graph neural networks, generative models, flow matching, and continuous-time neural models.\par
\vspace{4pt}
\textbf{Scientific computing:} ODE/PDE solvers, differentiable simulation, adjoint methods, inverse problems, data assimilation, uncertainty propagation, and numerical optimization.\par
\vspace{4pt}
\textbf{Programming:} Python, PyTorch, NumPy, SciPy, MATLAB, Git, and LaTeX.

\section*{Invited Talks}
\begin{itemize}
\item \textit{Topology-Aware Graph Neural Networks by Physics.} KGML Workshop: Leading the New Paradigm of AI for Science, University of Michigan, Ann Arbor, Aug. 2025.
\item \textit{Satellite trajectory modeling using machine learning methods.} International Symposium on Space Flight Dynamics, Beijing, Aug. 2022.
\item \textit{Orbit error propagation considering atmospheric density uncertainty.} Youth Innovation Promotion Association of the Chinese Academy of Sciences, Nanjing, Dec. 2021.
\end{itemize}

\section*{Professional Service}
\textbf{Reviewer:} ACM Multimedia (2025); IEEE ICDM (2025, 2026); AAAI, ICLR, KDD, and NeurIPS (2026).

\section*{Honors}
\textbf{Outstanding Graduate Student Award}, USTC, 2022 and 2023.\par
\vspace{3pt}
\textbf{First-Class Scholarship}, China University of Mining and Technology (Beijing), 2018 and 2019.
\end{document}
